\documentclass[../main]{subfiles}
\ifSubfilesClassLoaded{\setcounter{chapter}{1}}{}
\begin{document}

\chapter{Landasan Teori dan Konsep Dasar}

\begin{subcpmk}
  \item Sub-CPMK 1.1: Menjelaskan peran game engine dalam pengembangan game
\end{subcpmk}

% ============================================================
% MATERI POKOK
% ============================================================
\section{Pengertian Game dan Elemen-elemen Game}

\begin{learningoutcomes}
\item Mendefinisikan game secara formal dan akademis
\item Mengidentifikasi elemen-elemen pembentuk game (Mechanics, Dynamics, Aesthetics)
\item Menganalisis hubungan antara mechanics, dynamics, dan aesthetics
\item Menerapkan konsep MDA framework dalam analisis game
\end{learningoutcomes}

\subsection{Definisi Game}

\begin{definisi}[Game]
Game adalah aktivitas terstruktur yang melibatkan satu atau lebih pemain, dengan tujuan yang jelas, aturan yang mengatur interaksi, dan umpan balik yang memberikan informasi tentang kemajuan menuju tujuan \cite{Hunicke2004}.
\end{definisi}

Menurut \cite{Salen2004}, game memiliki karakteristik utama:
\begin{enumerate}
\item \textbf{Tujuan}: Setiap game memiliki tujuan yang harus dicapai pemain
\item \textbf{Aturan}: Aturan yang membatasi bagaimana pemain dapat mencapai tujuan
\item \textbf{Umpan Balik}: Informasi yang diberikan kepada pemain tentang status mereka
\item \textbf{Partisipasi Sukarela}: Pemain memilih untuk bermain dan menerima aturan
\end{enumerate}

\subsection{MDA Framework}

MDA (Mechanics, Dynamics, Aesthetics) merupakan framework analisis yang dikembangkan oleh \cite{Hunicke2004} untuk memahami game.

\subsubsection{Mechanics}

\begin{definisi}[Mechanics]
Mechanics adalah aturan dan sistem yang mengatur bagaimana game berfungsi. Mechanics mendefinisikan apa yang dapat dilakukan pemain dan bagaimana sistem game merespons aksi tersebut \cite{Hunicke2004}.
\end{definisi}

Contoh mechanics:
\begin{itemize}
\item Sistem pergerakan karakter
\item Mekanik pertempuran
\item Sistem inventori
\item Mekanik crafting
\item Sistem leveling dan progression
\end{itemize}

\subsubsection{Dynamics}

\begin{definisi}[Dynamics]
Dynamics adalah perilaku yang muncul dari interaksi mechanics saat game dimainkan. Dynamics adalah hasil dari mechanics yang bekerja bersama \cite{Hunicke2004}.
\end{definisi}

Contoh dynamics:
\begin{itemize}
\item Tension yang muncul dari time limit
\item Kompetisi antar pemain
\item Eksplorasi dan discovery
\item Risk vs reward decisions
\item Emergent gameplay
\end{itemize}

\subsubsection{Aesthetics}

\begin{definisi}[Aesthetics]
Aesthetics adalah pengalaman emosional yang dirasakan pemain saat bermain game. Aesthetics adalah respons emosional terhadap dynamics \cite{Hunicke2004}.
\end{definisi}

Jenis-jenis aesthetics menurut MDA framework:
\begin{enumerate}
\item \textbf{Sensation}: Game sebagai sense-pleasure
\item \textbf{Fantasy}: Game sebagai make-believe
\item \textbf{Narrative}: Game sebagai drama
\item \textbf{Challenge}: Game sebagai obstacle course
\item \textbf{Fellowship}: Game sebagai social framework
\item \textbf{Discovery}: Game sebagai uncharted territory
\item \textbf{Expression}: Game sebagai self-discovery
\item \textbf{Submission}: Game sebagai pastime
\end{enumerate}

\subsection{Hubungan Mechanics-Dynamics-Aesthetics}

\begin{figure}[h]
\centering
\begin{tikzpicture}[node distance=2cm]
    \node [block] (mechanics) {Mechanics\\Aturan Sistem};
    \node [block, below=of mechanics] (dynamics) {Dynamics\\Perilaku yang Muncul};
    \node [block, below=of dynamics] (aesthetics) {Aesthetics\\Pengalaman Emosional};
    
    \draw [line, ->] (mechanics) -- node[right] {Menghasilkan} (dynamics);
    \draw [line, ->] (dynamics) -- node[right] {Menciptakan} (aesthetics);
    \draw [line, ->, bend right=30] (aesthetics) to node[below] {Memengaruhi Desain} (mechanics);
\end{tikzpicture}
\caption{Hubungan Mechanics-Dynamics-Aesthetics}
\label{fig:mda_relationship}
\end{figure}

\subsection{Contoh Analisis Game Menggunakan MDA}

\begin{contoh}
\textbf{Analisis Game Platformer -- Super Mario Bros}

\textbf{Mechanics}: Jump mechanics dengan gravity, collision detection dengan platform, power-up system (mushroom, flower), enemy AI behavior, score system.

\textbf{Dynamics}: Timing jump untuk menghindari musuh, eksplorasi untuk menemukan power-up, risk-reward dalam mengambil coin, tension dari time limit.

\textbf{Aesthetics}: Challenge (mengatasi obstacle), Discovery (menemukan secret area), Sensation (visual dan audio feedback), Fantasy (menjadi karakter hero).
\end{contoh}

\begin{catatan}
MDA framework membantu desainer memahami bahwa perubahan pada mechanics akan mempengaruhi dynamics dan akhirnya aesthetics. Desain game adalah proses iteratif yang memerlukan testing dan refinement.
\end{catatan}

\section{Genre Game dan Karakteristiknya}

\begin{learningoutcomes}
\item Mengidentifikasi berbagai genre game utama
\item Menjelaskan karakteristik setiap genre
\item Menganalisis mekanik khas dari setiap genre
\item Menerapkan pengetahuan genre dalam desain game
\end{learningoutcomes}

\subsection{Pengertian Genre Game}

\begin{definisi}[Genre Game]
Genre game adalah kategori yang mengklasifikasikan game berdasarkan mekanik gameplay, setting, atau gaya permainan yang dominan \cite{Schell2019}.
\end{definisi}

Genre membantu pemain memahami apa yang diharapkan dari sebuah game dan membantu developer mengkomunikasikan konsep game mereka.

\subsection{Genre Game Utama}

\begin{figure}[ht]
    \centering
    \includegraphics[width=0.8\textwidth]{figures/game_genres.png}
    \caption{Visualisasi Ikon Genre Game Utama}
    \label{fig:game_genres}
\end{figure}

\subsubsection{Action Games}

\textbf{Karakteristik}: Fokus pada refleks dan koordinasi tangan-mata, gameplay cepat dan intens, sering melibatkan combat atau konflik langsung. Contoh: Call of Duty, Devil May Cry, Dark Souls.

\textbf{Sub-genre}: First-Person Shooter (FPS), Third-Person Shooter (TPS), Hack and Slash, Beat 'em up

\subsubsection{Adventure Games}

\textbf{Karakteristik}: Fokus pada narasi dan eksplorasi, puzzle-solving sebagai mekanik utama, interaksi dengan karakter dan lingkungan. Contoh: The Legend of Zelda, Uncharted, Life is Strange.

\subsubsection{Role-Playing Games (RPG)}

\textbf{Karakteristik}: Karakter development dan progression, sistem statistik (HP, MP, Level), quest dan mission-based gameplay. Contoh: Final Fantasy, The Elder Scrolls, The Witcher.

\subsubsection{Strategy Games}

\textbf{Karakteristik}: Fokus pada perencanaan dan taktik, resource management, decision-making jangka panjang. Contoh: Civilization, StarCraft, XCOM.

\subsubsection{Simulation Games}

\textbf{Karakteristik}: Simulasi sistem kehidupan nyata, realisme dan akurasi, open-ended gameplay. Contoh: The Sims, SimCity, Flight Simulator.

\subsubsection{Sports Games}

\textbf{Karakteristik}: Simulasi olahraga nyata, realisme dalam gameplay, multiplayer competitive. Contoh: FIFA, NBA 2K.

\subsubsection{Puzzle Games}

\textbf{Karakteristik}: Fokus pada problem-solving, logic-based challenges, sering tidak memiliki time pressure. Contoh: Tetris, Portal, The Witness.

\subsubsection{Platformer Games}

\textbf{Karakteristik}: Fokus pada jumping dan navigation, level design dengan platform, precision movement. Contoh: Super Mario Bros, Celeste, Hollow Knight.

\subsection{Hybrid Genre}

Banyak game modern menggabungkan elemen dari beberapa genre: Action-RPG, Survival Horror, Metroidvania.

\begin{catatan}
Genre bukanlah kategori yang kaku. Game modern sering menggabungkan elemen dari berbagai genre untuk menciptakan pengalaman yang unik.
\end{catatan}

\section{Perbandingan Game Engine}

Terdapat berbagai game engine populer. Pemilihan engine bergantung pada jenis game, tim, dan target platform. Berikut perbandingan singkat \cite{unityManual}:

\begin{table}[htbp]
\centering
\small
\begin{tabular}{|l|>{\raggedright\arraybackslash}p{3.2cm}|>{\raggedright\arraybackslash}p{3cm}|>{\raggedright\arraybackslash}p{2.2cm}|}
\hline
\textbf{Engine} & \textbf{Keunggulan} & \textbf{Kekurangan} & \textbf{Platform} \\
\hline
Unity & Mudah, gratis, multiplatform & Build size besar & PC, Mobile, Web \\
\hline
Unreal & Grafis tinggi, Blueprint & Kurva belajar curam & PC, Console \\
\hline
Godot & Open source, ringan & Ekosistem lebih kecil & PC, Mobile \\
\hline
\end{tabular}
\caption{Perbandingan Game Engine Populer}
\end{table}

\subsection{Unreal Engine}

Unreal Engine menggunakan C++ dan sistem Blueprint (visual scripting). Cocok untuk game AAA dengan grafis film. Kurva belajar cukup curam untuk pemula. Gratis dengan royalti setelah pendapatan tertentu.

\subsection{Godot}

Godot adalah engine open source dengan GDScript (mirip Python) atau C\#. Ringan, cocok untuk game 2D dan 3D sederhana. Komunitas lebih kecil dibanding Unity.

\subsection{Mengapa Unity untuk Buku Ini?}

Unity dipilih karena kemudahan untuk pemula dan kesesuaian dengan kurikulum:

\begin{enumerate}
  \item \textbf{Tutorial Resmi}: Unity Learn menyediakan Roll a Ball dan tutorial lain yang terstruktur langkah demi langkah
  \item \textbf{Asset Store}: Banyak aset gratis untuk pembelajaran
  \item \textbf{C\#}: Bahasa yang lebih mudah daripada C++ (Unreal), cocok untuk mahasiswa
  \item \textbf{Komunitas}: Forum dan dokumentasi sangat lengkap
  \item \textbf{Multiplatform}: Satu project bisa di-build untuk Windows, Mac, Android, iOS, Web
\end{enumerate}

\section{Lifecycle Pengembangan Game}

\begin{learningoutcomes}
\item Memahami fase-fase dalam lifecycle pengembangan game
\item Mengetahui aktivitas utama setiap fase
\item Memahami pentingnya planning dan documentation
\item Mampu mengidentifikasi milestone dalam development
\end{learningoutcomes}

\subsection{Pengertian Lifecycle Pengembangan Game}

\begin{definisi}[Game Development Lifecycle]
Lifecycle pengembangan game adalah proses sistematis dari konsep awal hingga rilis dan maintenance game, yang terdiri dari beberapa fase utama \cite{Gregory2018}.
\end{definisi}

\subsection{Fase-Fase Pengembangan Game}

\subsubsection{1. Pre-Production}

\textbf{Tujuan}: Membangun fondasi dan visi game

\textbf{Aktivitas Utama}: Brainstorming konsep, market research, prototyping core mechanics, menyusun GDD, budgeting dan resource planning, team assembly.

\textbf{Output}: GDD lengkap, prototype yang dapat dimainkan, art bible dan style guide, technical design document, project plan dan timeline.

\subsubsection{2. Production}

\textbf{Tujuan}: Membangun konten dan fitur game

\textbf{Aktivitas Utama}: Asset creation (art, audio, models), programming dan implementation, level design, integration dan testing, iterasi berdasarkan feedback.

\textbf{Fase dalam Production}:
\begin{enumerate}
\item \textbf{Alpha}: Core features implemented, masih banyak bug
\item \textbf{Beta}: Feature-complete, focus pada bug fixing dan polish
\item \textbf{Gold Master}: Final version siap untuk rilis
\end{enumerate}

\subsubsection{3. Post-Production}

\textbf{Tujuan}: Rilis, marketing, dan maintenance

\textbf{Aktivitas Utama}: Final testing dan QA, marketing dan promotion, launch preparation, distribution, post-launch support, updates dan patches, DLC development.

\subsection{Metodologi Pengembangan}

\textbf{Waterfall}: Pendekatan linear dengan fase berurutan. Keuntungan: terstruktur. Kekurangan: kurang fleksibel.

\textbf{Agile/Scrum}: Pendekatan iteratif dengan sprint pendek. Keuntungan: fleksibel, adaptif. Kekurangan: memerlukan komunikasi intensif.

\textbf{Hybrid}: Kombinasi waterfall untuk planning dan agile untuk execution.

\subsection{Milestone Penting}

Concept Approval, Prototype Complete, Pre-Alpha, Alpha, Beta, Gold, Launch.

\subsection{Best Practices}

Start Small, Prototype Early, Iterate Often, Document Everything, Test Continuously.


% ============================================================
% AKTIVITAS PEMBELAJARAN
% ============================================================
\begin{aktivitas}
  \item \textbf{Studi Literatur}: Pelajari game loop dan arsitektur game engine dari materi di atas dan Unity Documentation
  \item \textbf{Perbandingan}: Bandingkan Unity dengan satu game engine lain (Unreal/Godot). Buat tabel perbandingan singkat
  \item \textbf{Diskusi}: Diskusikan peran game engine dalam industri game development dengan teman
  \item \textbf{Praktek (opsional)}: Unity Learn tersedia untuk praktek mandiri---lihat Lampiran D
\end{aktivitas}

% ============================================================
% LATIHAN DAN REFLEKSI
% ============================================================
\begin{latihan}
  \item Jelaskan apa itu game loop dan mengapa penting!
  \item Sebutkan minimal 4 komponen yang disediakan game engine!
  \item Apa perbedaan GameObject dan Component dalam Unity?
  \item Mengapa Unity menggunakan C\# untuk scripting?
  \item \textbf{Refleksi}: Apa tantangan yang Anda rasakan dalam mempelajari game development?
\end{latihan}

% ============================================================
% ASESMEN
% ============================================================
\begin{asesmen}
\textbf{Instrumen Penilaian untuk Sub-CPMK 1.1}

\textbf{A. Pilihan Ganda}
\begin{enumerate}
  \item Peran game engine dalam pengembangan game adalah:
  \begin{enumerate}
    \item Hanya untuk rendering grafis
    \item Menyediakan komponen siap pakai (physics, audio, scripting)
    \item Menggantikan programmer sepenuhnya
    \item Hanya untuk game 2D
  \end{enumerate}
  
  \item Dalam Unity, GameObject minimal memiliki komponen:
  \begin{enumerate}
    \item Rigidbody
    \item Transform
    \item Collider
    \item Script
  \end{enumerate}
\end{enumerate}

\textbf{B. Essay}: Jelaskan alur game loop dan kaitannya dengan FPS (frame per second).
\end{asesmen}

% ============================================================
% CHECKLIST KOMPETENSI
% ============================================================
\begin{checklist}
  \item Saya dapat menjelaskan peran game engine dalam pengembangan game
  \item Saya memahami game loop (processInput, update, render) dan kaitannya dengan FPS
  \item Saya memahami arsitektur Entity-Component Unity
  \item Saya dapat menjelaskan layer arsitektur Unity (Native C++, C\# Layer, Editor)
\end{checklist}

% ============================================================
% RANGKUMAN
% ============================================================
\begin{rangkuman}
Bab ini membahas landasan teori pengembangan game, peran game engine, arsitektur Unity, dan perbandingan engine. Pemahaman ini menjadi fondasi untuk bab-bab praktik berikutnya.

\textbf{Kata Kunci}: \unity{Game Engine}, \unity{Game Loop}, \unity{GameObject}, \unity{Component}, \unity{MonoBehaviour}
\end{rangkuman}

\ifSubfilesClassLoaded{
  \renewcommand{\bibname}{Daftar Pustaka}
  \bibliographystyle{plain}
  \bibliography{../references}
}{}
\end{document}
